\section{The Foundational Role of the Paper in the Series}
\label{p11:sec:foundational}

\noindent
This paper is the programmatic and foundational text of the part of the series concerned with the praxis of flourishing, and the present section makes that role explicit, since a foundation is of use only if what it founds is stated. The section sets out the instruments the paper supplies to the subsequent papers, the questions it leaves open for them, and the interfaces at which they will connect to it. It is written for the reader of the series as a whole, and it may be passed over by the reader of this paper alone, for whom the preceding sections are complete in themselves.

\subsection{The instruments the paper supplies}
\label{p11:sec:found-instruments}

The paper supplies to the subsequent papers a set of instruments, each developed here in its general form and available there for application to particular cases.

The first is the \emph{lifecycle of a shared flourishing} as a real cycle (\xref{p11:sec:cycle}): the articulation of a shared life into the moments of cognition, practice, evaluation, reflection, and re-cognition, the thesis that the cycle is a self-sustaining movement whose persistence conditions the relation's existence, and the dynamical vocabulary of persistence, phase, and bifurcation in which the cycle's regimes are described. A subsequent paper concerned with any particular practice of a shared life may locate that practice within the cycle, asking at which moment it falls and how it bears on the cycle's persistence and phase.

The second is the \emph{operational concession} (\xref{p11:sec:bg-concession}): the methodological licence to adopt, where a practice must assess, a structured proxy for a value that is not symbolised, in the explicit awareness that the proxy is not the value. Every subsequent paper that assesses anything inherits this concession, and inherits with it the discipline of marking what is proxy and what is real, so that the series' results are read throughout as operationally warranted and not as measurements of the unsymbolised value.

The third is the \emph{skew of the return distribution} as a structured proxy for injustice (\xref{p11:sec:ev-skew}): the thesis that the injustice of a relation's value cycle has a statistical signature in the skewness of the distribution of returned value relative to a coupling-based baseline, and the associated apparatus of expected recyclability, coupling-based baseline, and the branching-process model of return. A subsequent paper concerned with the justice of any particular cycle of value, intimate or economic, may bring this measure to bear, asking whether the cycle returns to its generators or consigns one of them to the low-return tail.

The fourth is the \emph{distinction between the genuine and the forged cycle} (\xref{p11:sec:ref-blind}, \xref{p11:sec:reg-justice}): the thesis that a cycle may persist and present a high aggregate satisfaction while being, in its phase and its distribution, a catastrophic cycle wearing the appearance of a flourishing spiral, and that the orthogonal measures of skew and foreclosure distinguish the two. A subsequent paper concerned with any cycle that presents itself as good may apply this test, asking whether the appearance is borne out by the distribution beneath it.

The fifth is the \emph{cognitive and computational basis of appraisal} (\xref{p11:sec:appraisal}): the account of appraisal as continuous inference upon a partially observable state, conducted in unsymbolised value, with its decision-theoretic and free-energy formalisations. A subsequent paper concerned with how a relation registers its own state, or with how a practice is sustained or fails, may draw on this basis.

\subsection{The questions the paper leaves open}
\label{p11:sec:found-open}

The paper leaves open, for the subsequent papers and for other work, a set of questions it has posed without resolving. The branching-process model of return (\xref{p11:sec:ev-model}) is advanced in its form and not validated; the assignment of value to the acts of a shared life, on which any estimation of return depends, is identified as a substantial empirical problem and not solved; the formalisation of the lifecycle as a dynamical system is offered as a structural analogy under the standing qualification that the system's own law is reconfigured by what it generates; and the proxies for the unfolding of each party (\xref{p11:sec:ev-dynamics}) are acknowledged to reach least adequately the value that lies nearest the unsymbolised. These open questions are the paper's bequest to the work that follows, and their statement is part of the foundation, since a foundation that concealed its incompleteness would be the false ground the series has throughout refused.

\subsection{The interface to the study of particular cycles}
\label{p11:sec:found-interface}

The subsequent papers of this part of the series apply the instruments to particular cycles of value, and the present paper prepares their interface. A study of a particular cycle, on the apparatus developed here, proceeds by identifying the cycle's generators and their coupling, the value generated and its circulation, the distribution of return and its skew, and the cycle's persistence and phase; and it asks, of the cycle, whether it is a flourishing spiral, a catastrophic cycle, or a movement toward extinction, and whether its appearance of good is genuine or forged.

One such study is anticipated closely enough to be named. A consumption cycle, in which a good is desired, acquired, and enjoyed, and the desire renewed, may present every appearance of a flourishing cycle: the parties are satisfied, the cycle persists, and the satisfaction is real. The apparatus developed here permits the question whether such a cycle is a genuine flourishing spiral or a forged one, by directing attention past the satisfaction of the parties to the distribution of return across the whole network that sustains the cycle, including those whose labour produces the good consumed and who may be consigned, unseen by the satisfied parties, to the low-return tail. The skew of the return distribution, taken across the whole sustaining network and not only across the visible parties, is the instrument by which a consumption cycle that is internally satisfying may be found to be externally extractive, its internal holonomy positive and its full holonomy, reckoned across all whom it draws upon, negative. The study of such a cycle, as a case in which the genuine and the forged are distinguished by the apparatus of this paper, is the work of a subsequent paper, for which the present paper is the foundation, and to which it hands the instruments here developed.
